% Generated draft: requires mathematical and source review.
\documentclass[11pt]{article}
\usepackage[margin=25mm]{geometry}
\usepackage{fontspec}
\setmainfont{Latin Modern Roman}
\usepackage{amsmath,amssymb,mathtools,unicode-math}
\setmathfont{Latin Modern Math}
\usepackage{enumitem,needspace,xurl,hyperref}
\setlength{\parindent}{0pt}
\setlength{\parskip}{5pt}
\setlength{\emergencystretch}{3em}
\newcommand{\N}{\mathbb{N}}
\newcommand{\Z}{\mathbb{Z}}
\newcommand{\Q}{\mathbb{Q}}
\newcommand{\R}{\mathbb{R}}
\newcommand{\C}{\mathbb{C}}
\newcommand{\F}{\mathbb{F}}
\newcommand{\A}{\mathbb{A}}
\newcommand{\PP}{\mathbb{P}}
\newcommand{\E}{\mathbb{E}}
\newcommand{\eps}{\varepsilon}
\newcommand{\dd}{\,\mathrm{d}}
\begin{document}
\section*{UnsolvedMath ID: 600003}
\textbf{Source problem code:} catalogue entry
\subsection{Definitions and mathematical statement}
Equip $\R^2$ with the Lorentz bilinear form $g(u,v)=u_1v_1-u_2v_2$, so $\dd s^2=\dd x^2-\dd y^2$. A nonzero vector $v$ is space-like if $g(v,v)>0$, time-like if $g(v,v)<0$, and null if $g(v,v)=0$. Interpret an oval in the usual billiard setting as a smooth strictly convex closed curve $\gamma=\partial\Omega$ bounding a compact convex domain; the supplied question does not state its exact regularity assumptions.

At a point with non-null tangent, the Lorentz reflection in the tangent line sends an incoming vector $v$ to
\[
v-2\frac{g(v,n)}{g(n,n)}n,
\]
where $n$ is a nonzero Lorentz normal to the tangent. An ordinary periodic billiard trajectory has vertices $x_0,\ldots,x_{m-1}\in\gamma$, $m\ge2$, consecutive distinct vertices, and straight segments in $\overline\Omega$ satisfying this reflection law cyclically, with direction vectors taken up to positive rescaling. Call it space-like or time-like when its segment vectors have the respective sign of $g(v,v)$. These definitions describe the regular, non-null setting; they leave unspecified trajectories meeting a null tangent, null segments, or other singular collisions.

The source asks whether there is an analogue of Birkhoff's periodic-orbit existence theorem for every such Lorentz oval, with separate existence statements for space-like and time-like periodic trajectories. In particular, a precise analogue must specify which periods and, if relevant, which rotation or winding classes are to be prescribed, how many distinct orbits are guaranteed in each causal type, and which singularities are excluded or included. The source supplies no particular universal bound or fully quantified proposed theorem, so none is asserted here.

The second question asks for corresponding existence and multiplicity statements for billiards in bounded convex domains in pseudo-Euclidean spaces $\R^{p,q}$, with $p,q$ the signature parameters and bilinear form
\[
g_{p,q}(u,v)=\sum_{i=1}^{p}u_iv_i-\sum_{i=p+1}^{p+q}u_iv_i.
\]
The admissible domains, dimensions, periods, causal classes, and multiplicity targets for this extension are not specified in the record.

The supplied research summary reports: at period two, at least one space-like and one time-like orbit for Lorentz ovals and bounds of at least $p$ and $q$ diameters in the respective causal types in $\R^{p,q}$; a complete integrable picture for ellipses and ellipsoids, with Cayley-type criteria; and no stated general pseudo-Euclidean counterpart of the cited multidimensional Euclidean bounds. These are attributed background claims, not verified theorems proved by this draft or a general solution.
\subsection{Short English statement}
Does every Lorentzian convex billiard table have periodic trajectories of both causal types in the periods or winding classes expected from a suitable Birkhoff theorem, and what analogous guarantees hold in higher signature? The precise proposed bounds are missing. The supplied record labels the general question open with partial progress at period two and for ellipses or ellipsoids; its research summary is truncated and is not independent verification.
\subsection{Sources}
\begin{itemize}
\item Local source registry: \url{data/unsolvedmath/problems.json}, permanent UnsolvedMath ID $600003$, source code catalogue entry, including the supplied research summary. Dataset attribution: Ulam AI and UnsolvedMath Contributors; source attribution and rights are retained.
\item Tabachnikov, \emph{A Baker's Dozen of Problems} (2015), Section 3, unnumbered questions, as identified in the source list: \url{https://amj.math.stonybrook.edu/html-articles/Files-2015-2024/14-01/}. Source metadata records HTML access on July 29, 2026, and flags current status and rights for release review.
\item The supplied research summary attributes the period-two results to Khesin--Tabachnikov (2009), the ellipse and ellipsoid results to Dragovi\'c--Radnovi\'c (2012) and Adabrah--Dragovi\'c--Radnovi\'c (2019), and multidimensional Euclidean bounds to Farber--Tabachnikov (2002). It supplies no titles, precise theorem statements, or separate bibliographic links for these references.
\end{itemize}
\end{document}
